\documentclass[instructions.tex]{subfiles}
\begin{document}
 
\chapter{A la mort on est désabusé des vanités du monde.}
 
Ce n’est guère qu’en finissant la vie qu’on connait ce que c’est que la mort et ce que c’est que le monde; on comprend alors que la mort est le terme des vanités du siècle, et la fin de nos désordres ou de nos mérites. Voilà ce qui fait le regret ou la consolation des mourans.
 
\primo Oh! que la mort, en fermant les yeux du corps, ouvre clairement les yeux de l’âme! on voit pourquoi la vie nous était donnée, combien on a été aveugle de la passer dans l’oubli de Dieu; on connait que tout, en ce monde, n’est que néant; hors aimer Dieu et le servir.
 
L’empereur Septime Sévère disait en mourant : J’ai possédé le plus grand empire de l’univers, j’ai été tout ce que peut être un mortel; il ne me reste rien qui puisse me servir dans l’état où je suis. Un pape disait quelques momens avant sa mort : J’ai eu les clefs du Ciel entre les mains; mais il vaudrait mieux pour moi avoir eu celles d’un pauvre monastère.
 
Quand tout est passé, c’est alors qu’on se repent de n’avoir pas fait ce qu’on devait faire; mais qu’il est bien tard de regretter le passé, quand on est hors d’état d’y remédier! Qu’il est fâcheux de comprendre ses égaremens, quand on ne peut plus revenir sur ses pas; et de voir qu’on a mal vécu, quand on n’a plus de temps à vivre!
 
Pour éviter tous ces repentirs affligeans, pensez à présent des choses du monde comme vous en penserez à l’heure de la mort. Que vous revient-il de tout ce qui fait aujourd’hui le sujet de votre confusion? disait saint Paul aux Romains. Je vous fais la même demande : Que vous reste-t-il de toutes les choses agréables et amusantes qui vous ont occupé? Que sont devenus ces spectacles, ces fréquentations, ces débauches, et tout ce qui vous a charmé dans les biens et les plaisirs de la vie? Il ne vous en reste qu’un souvenir, avec le regret d’y avoir perdu ou risqué votre âme. Toutes ces choses ne changeront pas avec le temps; autant de fois que vous vous y livrerez, vous n’en retirerez jamais d’autres fruits que des remords et des repentirs. Voilà tout ce qui vous restera au moment de la mort.
 
Quand vous auriez joui jusqu’à présent de tous les bonheurs et de toutes les délices de la terre, ou quand vous auriez éprouvé toutes les disgrâces, et tout ce qu’on peut souffrir, un moment après, tout cela ne sera plus. Il ne vous restera des biens et des plaisirs dont vous aurez joui, que le chagrin de les perdre; et des maux que vous aurez soufferts, que le repentir d’avoir abusé de ce qui pouvait vous procurer le ciel; votre regret sera de vous être servi, pour votre damnation, de ce qui devait vous mériter une gloire éternelle : regret accablant, puisque l’éternité ne pourra jamais en adoucir l’amertume. J’ai vécu vingt années sur le trône, disait en mourant un grand roi, et il ne m’en reste rien! Plût à Dieu que ces vingt années de règne eussent été vingt années de retraite, employées pour le service de Dieu et pour mon salut!
 
\secundo La mort met tous les hommes au même niveau; tout est alors également passé pour les uns et pour les autres, mais bien différemment pour les gens de bien et pour les pécheurs. Voyez un saint Antoine mourant à l’âge de cent et cinq ans; qui pourra comprendre la joie de son âme à la vue du temps et de l’éternité : du temps, avec lequel viennent de passer toutes les rigueurs de sa vie pénitente; de l’éternité, avec laquelle vont commencer tous les plaisirs que procure la jouissance de Dieu?
 
Représentez-vous, au contraire, le malheureux Caïn; quand il aurait vécu 5000 ans dans les délices et dans l’abondance de toutes choses, que lui en resterait-il à présent? On est bien aveugle, dit un Père, de se laisser surprendre par l’amorce des plaisirs et des biens de la vie, sans considérer combien la fin en est amère.
 
L’impératrice Isabelle étant morte à Tolède, on fit transporter son corps à Grenade. A l’ouverture du cercueil, le visage de cette princesse parut si défiguré et si hideux, que saint François de Borgia, duc de Gandie, présent à cette cérémonie, s’écria : Est-ce donc là cette reine qui avait tant de charmes, il n’y a que quelques jours? Quoi! voilà ce visage autrefois si brillant! O mon Dieu! que l’homme est misérable d’attacher son cœur à des objets si frivoles! et que celles qui ont tant de soin d’une beauté si fragile sont aveugles! Il passa la nuit suivante à déplorer les folies du monde, se fit ensuite religieux dans la compagnie de Jésus, et s’y sanctifia.
 
On ne goûte point de plus douce consolation à la mort, que par la pensée d’avoir vécu dans le mépris des plaisirs et des grandeurs, et d’avoir servi Dieu.
 
\section{Résolutions}
 
\primo C’en est fait : je vais quitter, dès ce moment, tout ce qui pourra me causer des regrets à la vue de la mort.--
 
\secundo Je renonce donc aux vains fantômes dont le penchant que j’ai à l’orgueil, aux richesses, aux plaisirs dangereux, me remplit l’imagination et le cœur, et je vais désormais m’attacher uniquement à ce qui me donnera des consolations à la mort.

\prayer{O mon Dieu! faites que j’embrasse tout de bon, dès à présent, et pour le reste de ma vie, la pénitence et la pratique soigneuse de tous mes devoirs.}
 
\end{document}
